\documentclass[11pt,a4paper]{article}

% ============================================================
% PACKAGES
% ============================================================
\usepackage[utf8]{inputenc}
\usepackage[T1]{fontenc}
\usepackage[english]{babel}
\usepackage{amsmath, amssymb, amsthm}
\usepackage{geometry}
\geometry{margin=2.5cm}
\usepackage{graphicx}
\usepackage{xcolor}
\usepackage{listings}
\usepackage{hyperref}
\usepackage{booktabs}
\usepackage{array}
\usepackage{tikz}
\usetikzlibrary{arrows.meta, positioning, shapes, calc}
\usepackage{enumitem}

% ============================================================
% CONFIGURATION
% ============================================================
\hypersetup{
    colorlinks=true,
    linkcolor=blue,
    urlcolor=blue,
    citecolor=blue
}

\lstset{
    language=Python,
    basicstyle=\ttfamily\small,
    keywordstyle=\color{blue}\bfseries,
    commentstyle=\color{gray}\itshape,
    stringstyle=\color{red!70!black},
    numbers=left,
    numberstyle=\tiny\color{gray},
    frame=single,
    breaklines=true,
    showstringspaces=false,
    tabsize=4,
    literate={≤}{{$\leq$}}1 {≥}{{$\geq$}}1 {≠}{{$\neq$}}1
}

\newcommand{\bigO}{\mathcal{O}}
\newcommand{\bigOmega}{\Omega}
\newcommand{\bigTheta}{\Theta}
\newcommand{\code}[1]{\texttt{\detokenize{#1}}}

% Difficulty stars
\newcommand{\easy}{$\star$}
\newcommand{\medium}{$\star\star$}
\newcommand{\hard}{$\star\star\star$}

% ============================================================
% TITLE
% ============================================================
\title{\textbf{Algorithms: Complexity and Optimization}\\
\large Exercise Book}
\author{}
\date{}

\begin{document}

\maketitle

\tableofcontents
\newpage

% ============================================================
% INTRODUCTION
% ============================================================
\section*{Introduction}

This exercise book accompanies the course \emph{Algorithms: Complexity and Optimization}. Each section follows the structure of the course and contains two types of exercises:
\begin{itemize}
    \item \textbf{Theoretical exercises} (tracing, proving, drawing) --- to be solved on paper first.
    \item \textbf{Programming exercises} (coding, testing) --- to be implemented in Python.
\end{itemize}

Each exercise is graded by difficulty:
\begin{itemize}
    \item \easy\ --- warm-up, direct application of the course.
    \item \medium\ --- requires some reasoning.
    \item \hard\ --- challenging, requires proof or deep understanding.
\end{itemize}

\newpage

% ============================================================
% 1. ASYMPTOTIC NOTATION
% ============================================================
\section{Asymptotic Notation}

\subsection{Theoretical Exercises}

\begin{enumerate}[label=\textbf{1.\arabic*}]

    \item[\easy] Simplify $O(4n^3 + 2n^2 + 100)$. Justify your answer.

    \item[\easy] True or false: $2^n = O(n^2)$? Justify.

    \item[\easy] Find the tightest $\Theta$ for $5n \log n + 3n + 2$.

    \item[\medium] Prove that $n^2 + 3n = \Omega(n^2)$.

    \item[\medium] Order the following functions by growth rate:
    \[
    \sqrt{n}, \quad n, \quad n \log n, \quad \log n, \quad n^2, \quad 2^n.
    \]

    \item[\medium] Prove that $n \log n = O(n^2)$ but $n^2 \neq O(n \log n)$.

    \item[\medium] Give an example of an algorithm with $O(1)$ time complexity but $O(n)$ space complexity. Justify.

    \item[\medium] Prove that $3n^2 + 5n + 7 = O(n^2)$ by giving explicit constants $c$ and $n_0$.

\end{enumerate}

% ============================================================
% 2. DATA STRUCTURES
% ============================================================
\section{Data Structures}

\subsection{Theoretical Exercises}

\begin{enumerate}[label=\textbf{2.\arabic*}]

    \item[\easy] What is the complexity of accessing the last element of an array? The first element of a linked list? Justify.

    \item[\easy] What is the complexity of inserting an element at the beginning of an array? At the end? In the middle? Justify.

    \item[\medium] Implement a queue using two stacks. What is the complexity of each operation (enqueue, dequeue)?

    \item[\medium] Implement a stack using two queues. What is the complexity of each operation (push, pop)?

    \item[\easy] Draw the binary search tree (BST) obtained by inserting the following keys in this order:
    \[
    50, 30, 70, 20, 40, 60, 80.
    \]
    What is the height of the tree?

    \item[\easy] What is the height of a BST with $n$ nodes in the best case? In the worst case? Justify.

    \item[\medium] Represent a graph with 5 vertices and 6 edges as an adjacency matrix and as an adjacency list. Compare the space used by the two representations.

    \item[\hard] Prove that the height of a balanced BST with $n$ nodes is $O(\log n)$.

    \item[\medium] Why does a degenerate BST have $O(n)$ search complexity? Give an insertion order that produces a degenerate BST.

    \item[\medium] Compare the space complexity of an adjacency matrix and an adjacency list for a sparse graph. Which is better, and why?

\end{enumerate}

\subsection{Programming Exercises}

\begin{enumerate}[label=\textbf{2.\arabic*}, start=11]

    \item[\easy] Write a function \code{find_max(arr)} that returns the maximum element of an array. Test it on \code{[3, 1, 4, 1, 5, 9, 2, 6]}.
    \begin{lstlisting}
def find_max(arr):
    # your code here
    pass

# Test
print(find_max([3, 1, 4, 1, 5, 9, 2, 6]))  # Expected: 9
    \end{lstlisting}

    \item[\easy] Write a function \code{reverse_array(arr)} that reverses an array \emph{in place}, in $O(n)$ time and $O(1)$ space. Test it on \code{[1, 2, 3, 4, 5]}.
    \begin{lstlisting}
def reverse_array(arr):
    # your code here
    pass

# Test
arr = [1, 2, 3, 4, 5]
reverse_array(arr)
print(arr)  # Expected: [5, 4, 3, 2, 1]
    \end{lstlisting}

    \item[\easy] Write a function \code{is_sorted(arr)} that returns \code{True} if the array is sorted in ascending order, \code{False} otherwise. Test it on \code{[1, 2, 3]} and \code{[1, 3, 2]}.
    \begin{lstlisting}
def is_sorted(arr):
    # your code here
    pass

# Test
print(is_sorted([1, 2, 3]))  # Expected: True
print(is_sorted([1, 3, 2]))  # Expected: False
    \end{lstlisting}

    \item[\medium] Implement a linked list with the following operations: \code{insert_at_head}, \code{insert_at_tail}, \code{delete}, \code{search}, \code{print_list}. Test it on a sequence of 10 operations.
    \begin{lstlisting}
class Node:
    def __init__(self, value):
        self.value = value
        self.next = None

class LinkedList:
    def __init__(self):
        self.head = None

    def insert_at_head(self, value):
        # your code here
        pass

    def insert_at_tail(self, value):
        # your code here
        pass

    def delete(self, value):
        # your code here
        pass

    def search(self, value):
        # your code here
        pass

    def print_list(self):
        # your code here
        pass
    \end{lstlisting}

    \item[\medium] Implement a stack using a Python list. Provide \code{push}, \code{pop}, \code{peek}, and \code{is_empty}. Test it.
    \begin{lstlisting}
class Stack:
    def __init__(self):
        self.items = []

    def push(self, x):
        # your code here
        pass

    def pop(self):
        # your code here
        pass

    def peek(self):
        # your code here
        pass

    def is_empty(self):
        # your code here
        pass
    \end{lstlisting}

    \item[\medium] Implement a queue using \code{collections.deque}. Provide \code{enqueue}, \code{dequeue}, \code{peek}, and \code{is_empty}. Test it.
    \begin{lstlisting}
from collections import deque

class Queue:
    def __init__(self):
        self.items = deque()

    def enqueue(self, x):
        # your code here
        pass

    def dequeue(self):
        # your code here
        pass

    def peek(self):
        # your code here
        pass

    def is_empty(self):
        # your code here
        pass
    \end{lstlisting}

    \item[\medium] Implement a queue using two stacks. What is the amortized complexity of \code{dequeue}? Test it on 10 enqueues and 5 dequeues.
    \begin{lstlisting}
class QueueTwoStacks:
    def __init__(self):
        self.inbox = []
        self.outbox = []

    def enqueue(self, x):
        # your code here
        pass

    def dequeue(self):
        # your code here
        pass
    \end{lstlisting}

    \item[\hard] Write a function \code{is_balanced(s)} that checks whether a string of parentheses is balanced. Test it on \code{"(())()"}, \code{"(()"}, and \code{")("}.
    \begin{lstlisting}
def is_balanced(s):
    # your code here
    pass

# Test
print(is_balanced("(())()"))  # Expected: True
print(is_balanced("(()"))     # Expected: False
print(is_balanced(")("))      # Expected: False
    \end{lstlisting}

\end{enumerate}

% ============================================================
% 3. SEARCH IN BST AND SORTED LIST
% ============================================================
\section{Search in BST and Sorted List}

\subsection{Theoretical Exercises}

\begin{enumerate}[label=\textbf{3.\arabic*}]

    \item[\easy] Trace the iterative dichotomy search on the array
    \[
    [1, 3, 5, 7, 9, 11, 13]
    \]
    for the value $7$. How many comparisons are performed?

    \item[\easy] What is the complexity of searching in a degenerate BST? In a balanced BST? Justify.

    \item[\medium] Write the iterative and recursive versions of dichotomy. Compare their space complexity.

    \item[\medium] Why is the recursive version of dichotomy a tail call? What does this imply in a language that supports tail-call optimization?

    \item[\hard] Prove that any comparison-based search in a sorted array of $n$ elements requires $\Omega(\log n)$ comparisons in the worst case.

    \item[\medium] Write a function that finds the \textbf{first} element $\geq x$ in a sorted array, in $O(\log n)$ time.

    \item[\medium] Write a function that finds the \textbf{number of occurrences} of $x$ in a sorted array, in $O(\log n)$ time.

    \item[\hard] Write a function that finds the \textbf{$k$-th smallest} element in a BST.

\end{enumerate}

% ============================================================
% 4. SORTING ALGORITHMS
% ============================================================
\section{Sorting Algorithms}

\subsection{Theoretical Exercises}

\begin{enumerate}[label=\textbf{4.\arabic*}]

    \item[\easy] Trace bubble sort on the array
    \[
    [5, 2, 4, 6, 1, 3].
    \]
    Count the number of swaps.

    \item[\easy] Trace merge sort on the array
    \[
    [38, 27, 43, 3, 9, 82, 10].
    \]
    Draw the recursion tree and the merging steps.

    \item[\easy] Trace quick sort (with the last element as pivot) on the array
    \[
    [10, 7, 8, 9, 1, 5].
    \]
    Show the partition at each step.

    \item[\medium] Why does quick sort have an $O(n^2)$ worst case? Give an example of input that triggers it, and explain how to avoid it with a randomized pivot.

    \item[\medium] Prove that merge sort is $O(n \log n)$ in all cases. Write the recurrence and solve it.

    \item[\hard] Prove that any comparison-based sorting algorithm requires $\Omega(n \log n)$ comparisons in the worst case.

\end{enumerate}

\subsection{Programming Exercises}

\begin{enumerate}[label=\textbf{4.\arabic*}, start=7]

    \item[\easy] Implement insertion sort in Python and test it on 5 examples (empty array, one element, sorted array, reverse-sorted array, random array).

    \item[\medium] Implement merge sort in Python and test it on \code{[38, 27, 43, 3, 9, 82, 10]}.

    \item[\medium] Implement quick sort with a randomized pivot. Test it on a sorted array, a reverse-sorted array, and a random array.

    \item[\medium] Compare experimentally bubble sort, quick sort, and merge sort on arrays of size $n = 1000$, $n = 10{,}000$, and $n = 100{,}000$. Report the execution times and discuss the results.

\end{enumerate}

% ============================================================
% 5. GREEDY ALGORITHMS
% ============================================================
\section{Greedy Algorithms}

\subsection{Theoretical Exercises}

\begin{enumerate}[label=\textbf{5.\arabic*}]

    \item[\easy] Apply the greedy coin change algorithm to $S = 87$ with the European system $\{1, 2, 5, 10, 20, 50\}$. Give the dictionary of coins used and the total number of coins.

    \item[\easy] Show that the greedy algorithm fails for the coin system $\{1, 3, 4\}$ with $S = 6$. What is the optimal change? How many coins does the greedy algorithm use?

    \item[\medium] Generate all possible changes for $S = 5$ with coins $\{1, 2, 3\}$, without considering the order. How many are there?

    \item[\medium] Draw the tree of all possible changes for $S = 4$ with coins $\{1, 2\}$. Highlight the valid changes (leaves where the remaining sum is 0).

    \item[\medium] Find another coin system (different from $\{1, 3, 4\}$) where the greedy algorithm fails. Give the counter-example and the optimal change.

\end{enumerate}

\subsection{Programming Exercises}

\begin{enumerate}[label=\textbf{5.\arabic*}, start=6]

    \item[\medium] Implement \code{interval_scheduling(tasks)} that returns the maximum number of non-overlapping tasks. Test it on a small example.
    \begin{lstlisting}
def interval_scheduling(tasks):
    # tasks = [(start, end), ...]
    # your code here
    pass

# Test
tasks = [(1, 3), (2, 5), (3, 4), (5, 7), (6, 8)]
print(interval_scheduling(tasks))  # Expected: 3
    \end{lstlisting}

\end{enumerate}

% ============================================================
% 6. DYNAMIC PROGRAMMING
% ============================================================
\section{Dynamic Programming}

\subsection{Theoretical Exercises}

\begin{enumerate}[label=\textbf{6.\arabic*}]

    \item[\easy] Fill the DP table for the coin change problem with $S = 10$ and coins $\{1, 3, 4\}$. Give the final answer.

    \item[\medium] Draw the tree of balanced parentheses for $n = 4$. How many leaves (valid strings) are there?

    \item[\medium] Write the memoized and tabulated versions of Fibonacci. Test both on $n = 30$. Compare their space complexity and discuss the practical differences.

    \item[\medium] Write the DP for the generalized coin change problem (arbitrary coin system). Test it on $\{1, 3, 4\}$ with $S = 6$, and compare with the greedy result.

    \item[\medium] Write the DP for the Longest Common Subsequence (LCS). Test it on the strings \code{"ABCBDAB"} and \code{"BDCABA"}.
    
\end{enumerate}

% ============================================================
% 7. GRAPHS
% ============================================================
\section{Graphs}

\subsection{Programming Exercises}

\begin{enumerate}[label=\textbf{7.\arabic*}]

    \item[\medium] Implement a graph as an adjacency list. Provide \code{add_edge(u, v)} and \code{print_graph()}. Test it on a graph with 5 vertices.
    \begin{lstlisting}
class Graph:
    def __init__(self):
        self.adj = {}

    def add_edge(self, u, v):
        # your code here
        pass

    def print_graph(self):
        # your code here
        pass
    \end{lstlisting}

    \item[\medium] Implement BFS on a graph represented as an adjacency list. Test it on a small graph and print the order of visits.
    \begin{lstlisting}
from collections import deque

def bfs(graph, start):
    # your code here
    pass
    \end{lstlisting}

    \item[\medium] Implement DFS (recursive and iterative) on the same graph. Compare the order of visits with BFS.
    \begin{lstlisting}
def dfs_recursive(graph, start):
    # your code here
    pass

def dfs_iterative(graph, start):
    # your code here
    pass
    \end{lstlisting}

    \item[\hard] Implement Dijkstra's algorithm for shortest paths in a weighted graph. Test it on a graph with 5 vertices and positive weights.
    \begin{lstlisting}
import heapq

def dijkstra(graph, source):
    # your code here
    pass
    \end{lstlisting}

\end{enumerate}

\end{document}